\chapter{Leg ulcers}

\section*{Definition and Overview}
Leg ulcers are chronic, slow-healing open sores or wounds that typically develop on the lower extremities, particularly around the ankle or lower calf. They are characterised by a loss of epidermis and varying degrees of dermis and subcutaneous tissue disruption. Leg ulcers are primarily caused by underlying vascular disease, with venous insufficiency being the most common, followed by arterial insufficiency and neuropathic causes. These lesions represent a significant clinical challenge due to their chronic nature, high rate of recurrence, susceptibility to secondary infection, and profound negative impact on the patient's quality of life.

\section*{Epidemiology}
Leg ulcers are a common condition, particularly among the elderly and individuals with chronic medical comorbidities. They affect approximately 1\% to 2\% of the adult population globally, with the prevalence rising sharply to over 3\% to 5\% in individuals aged over 80 years. In Australia, leg ulcers represent a major healthcare burden, accounting for significant health system expenditure and nursing resources, particularly in community and aged-care settings. Women are slightly more affected than men, primarily due to higher rates of chronic venous insufficiency and longer life expectancy.

\section*{Aetiology and Pathophysiology}
The development of leg ulcers is driven by vascular dysfunction, neurological impairment, or systemic diseases that impair local tissue perfusion and healing.
\begin{itemize}
    \item \textbf{Venous insufficiency:} Valvular incompetence in the deep or superficial veins leads to venous hypertension, blood pooling, extravasation of macromolecules, and localised tissue ischaemia.
    \item \textbf{Arterial insufficiency:} Atherosclerotic narrowing or occlusion of lower limb arteries impairs oxygen and nutrient delivery to tissues, leading to ischaemic necrosis.
    \item \textbf{Neuropathy:} Sensory loss (often due to diabetes) leads to unnoticed trauma, repetitive pressure, and subsequent ulceration, particularly on the plantar aspect of the foot.
    \item \textbf{Diabetes mellitus:} A combination of microvascular disease, macrovascular disease, and neuropathy severely impairs tissue healing and increases infection susceptibility.
    \item \textbf{Vasculitis:} Inflammatory conditions affecting blood vessels can cause local thrombosis, tissue necrosis, and ulceration.
    \item \textbf{Trauma:} Minor physical injury or friction in a limb with pre-existing compromised circulation can trigger non-healing wounds.
\end{itemize}

\section*{Clinical Features}
The clinical presentation varies depending on the underlying cause, with characteristic differences in location, appearance, and pain.
\begin{itemize}
    \item \textbf{Venous ulcers:} Typically located on the medial gaiter region of the lower leg; large, shallow, with irregular borders, sloping edges, and frequent exudate, associated with hyperpigmentation (haemosiderin staining), varicose eczema, and oedema.
    \item \textbf{Arterial ulcers:} Commonly found on distal points (toes, heels, bony prominences); punched-out appearance with well-defined borders, a pale or necrotic wound bed, minimal exudate, and severe pain that worsens with leg elevation.
    \item \textbf{Neuropathic ulcers:} Typically occur on pressure-bearing areas of the foot; surrounded by calloused skin, painless, and deep, often leading to underlying osteomyelitis.
    \item \textbf{Local pain and tenderness:} Varies from dull, aching discomfort in venous disease to sharp, severe, claudicating pain in arterial disease.
    \item \textbf{Surrounding skin changes:} Shiny, hairless, cool skin in arterial disease; dry, itchy, indurated, or hyperpigmented skin (lipodermatosclerosis) in venous disease.
\end{itemize}

\section*{Differential Diagnosis}
Differentiating the exact cause of a lower limb ulcer is crucial for directing appropriate treatment, as incorrect therapies can lead to severe complications.
\begin{itemize}
    \item \textbf{Malignancy:} Squamous cell carcinoma (Marjolin's ulcer), basal cell carcinoma, or melanoma presenting as a non-healing atypical ulcer.
    \item \textbf{Pyoderma gangrenosum:} A painful, rapidly progressing neutrophilic dermatosis often associated with inflammatory bowel disease or rheumatoid arthritis.
    \item \textbf{Infectious ulcers:} Ulcers caused by atypical infections such as tuberculosis, atypical mycobacteria, deep fungal infections, or cutaneous leishmaniasis.
    \item \textbf{Sickle cell anaemia:} Chronic ulceration secondary to microvascular occlusion and tissue hypoxia in patients with sickle cell disease.
    \item \textbf{Calciphylaxis:} A rare, painful condition involving vascular calcification and skin necrosis, typically in patients with end-stage renal failure.
\end{itemize}

\section*{When to Seek Medical Care}
Prompt medical evaluation is necessary for any new or non-healing leg wound to prevent severe tissue damage or systemic infection.
\textbf{Call 000 (or local emergency services) for:}
\begin{itemize}
    \item Sudden, heavy bleeding from the ulcer or varicose veins in the leg.
    \item Signs of systemic sepsis, including high fever, rigors, rapid breathing, and confusion.
    \item A rapidly spreading red, hot, and extremely painful rash on the leg (severe cellulitis).
    \item Sudden onset of a cold, pale, and severely painful leg indicating acute limb-threatening arterial occlusion.
\end{itemize}
Other reasons to seek medical care include noticing a new leg sore that does not heal within two weeks, an increase in pain, foul-smelling discharge, or swelling in the lower limb.

\section*{Diagnosis and Investigations}
A combination of clinical assessment, vascular imaging, and laboratory tests is required to establish the primary aetiology of the leg ulcer.
\begin{itemize}
    \item \textbf{Ankle-Brachial Pressure Index (ABPI):} A non-invasive test measuring the ratio of systolic blood pressure in the ankle to that in the arm, critical for ruling out arterial disease before applying compression therapy.
    \item \textbf{Duplex ultrasound:} Visualises venous anatomy and flow to identify superficial or deep venous reflux, obstruction, or deep vein thrombosis (DVT).
    \item \textbf{Wound swab or culture:} Performed if clinical signs of infection are present (e.g., increased purulent exudate, worsening pain, erythema) to identify bacterial pathogens.
    \item \textbf{Skin biopsy:} Indicated for atypical ulcers, those failing to improve after 12 weeks of standard therapy, or if malignancy or vasculitis is suspected.
    \item \textbf{Blood tests:} Blood glucose, HbA1c, and inflammatory markers (CRP, ESR) to screen for diabetes, systemic inflammation, or nutritional deficiencies.
\end{itemize}

\section*{Management}
The management of leg ulcers is multi-disciplinary and depends on the specific primary cause, focusing on wound care, compression, and addressing underlying vascular disease.
\begin{itemize}
    \item \textbf{Compression therapy:} The cornerstone of venous ulcer management, utilising multi-layer compression bandages or compression stockings to improve venous return and reduce oedema (contraindicated if ABPI is less than 0.8).
    \item \textbf{Wound debridement and dressings:} Surgical, mechanical, or autolytic removal of devitalised tissue, combined with specialised dressings (e.g., alginates, hydrocolloids, foams) to maintain a moist healing environment.
    \item \textbf{Revascularisation:} Surgical or endovascular procedures (e.g., angioplasty, arterial bypass, or venous ablation) to restore perfusion in arterial disease or correct venous reflux.
    \item \textbf{Infection control:} Appropriate use of systemic antibiotics if cellulitis or systemic infection is present, as topical antibiotics are generally avoided to prevent sensitisation.
    \item \textbf{Lifestyle and nutritional support:} Optimising diet (adequate protein, zinc, and vitamins A and C), managing body weight, controlling blood glucose levels, and promoting gentle exercise to activate the calf muscle pump.
\end{itemize}

\section*{Complications}
\begin{itemize}
    \item \textbf{Cellulitis and erysipelas:} Acute bacterial infections of the skin and subcutaneous tissues surrounding the ulcer.
    \item \textbf{Osteomyelitis:} Spread of infection to the underlying bone, particularly common in deep neuropathic or diabetic foot ulcers.
    \item \textbf{Sepsis:} Life-threatening systemic inflammatory response triggered by a severe wound infection.
    \item \textbf{Malignant transformation:} Development of squamous cell carcinoma (Marjolin's ulcer) in long-standing, chronic wounds.
    \item \textbf{Contact dermatitis:} Allergic reactions to topical dressings, ointments, or preservatives used in wound care.
\end{itemize}

\section*{Prognosis and Outcomes}
The prognosis for leg ulcers varies significantly based on the underlying aetiology, patient compliance with treatment, and comorbidities. Venous ulcers generally have a high healing rate (up to 70\% to 80\% within six months) with appropriate compression therapy, but recurrence rates are high (up to 70\% within five years if compression is discontinued). Arterial and diabetic ulcers have a more guarded prognosis, requiring surgical revascularisation or intensive multi-disciplinary care to prevent tissue loss or lower limb amputation.

\section*{Prevention and Self-Care}
\begin{itemize}
    \item Wear compression stockings daily as prescribed to prevent the recurrence of venous ulcers.
    \item Inspect legs and feet daily for any scratches, cuts, redness, or blisters, especially if diabetic.
    \item Keep the skin clean and moisturised to prevent drying and cracking, taking care to dry thoroughly between the toes.
    \item Elevate the legs above the level of the heart when resting to reduce venous pressure and swelling.
    \item Engage in regular physical activity, such as walking, to maintain joint mobility and improve leg circulation.
\end{itemize}

\section*{Sources and Further Reading}
The following reputable organisations provide further information and clinical guidelines on leg ulcers and wound care:
\begin{itemize}
    \item Wounds Australia. \textit{Clinical practice guidelines for the management of venous leg ulcers.} https://www.woundsaustralia.com.au/
    \item NHS. \textit{Overview of venous leg ulcers.} https://www.nhs.uk/
    \item Healthdirect Australia. \textit{Leg ulcers information and care.} https://www.healthdirect.gov.au/
    \item Mayo Clinic. \textit{Understanding arterial and venous ulcers.} https://www.mayoclinic.org/
    \item National Institute for Health and Care Excellence (NICE). \textit{Leg ulcer guidelines and pathways.} https://www.nice.org.uk/
\end{itemize}
